\documentclass[12pt,oneside]{report}

\usepackage[a4paper,margin=1in,bindingoffset=0.25in]{geometry}
\usepackage{amsmath,amssymb}
\usepackage{booktabs}
\usepackage{longtable}
\usepackage{array}
\usepackage{microtype}
\usepackage{setspace}
\usepackage[hidelinks]{hyperref}
\usepackage{xcolor}

\onehalfspacing
\setlength{\parindent}{0.3in}
\setlength{\parskip}{0pt}
\setlength{\emergencystretch}{2em}
\renewcommand{\arraystretch}{1.15}

\newcommand{\systemname}{GEODE}
\newcommand{\ba}{\operatorname{BA}}
\newcommand{\utility}{U}
\newcommand{\thesistitle}{Adaptation Utility and Transactional Safety in
Open-World Recognition Systems}

\begin{document}
\hypersetup{pageanchor=false}

\begin{titlepage}
\centering
\vspace*{0.8in}

{\Large\bfseries \thesistitle\par}
\vspace{0.3in}
{\large A Study of Why Qualified Recognition, Discovery, and Adaptation
Stages Fail to Compose\par}

\vspace{0.8in}

{\large by\par}
\vspace{0.15in}
{\Large Moazzam Abdullah Khan\par}

\vfill

A thesis-style report submitted in partial fulfillment\\
of the requirements for the degree of\\[0.15in]
{\large Master of Science\par}

\vspace{0.5in}

Department and University Name\\
City, Country

\vspace{0.5in}

July 2026
\end{titlepage}

\hypersetup{pageanchor=true}
\pagenumbering{roman}

\chapter*{Approval}
\addcontentsline{toc}{chapter}{Approval}

This thesis-style report has been prepared for academic review. Institutional
approval language, degree title, department name, supervisor names, signatures,
and submission date must be replaced with the requirements of the submitting
university before formal deposit.

\vspace{0.6in}

\begin{tabular}{p{2.3in}p{2.3in}}
\hrulefill & \hrulefill\\
Supervisor & Date\\[0.5in]
\hrulefill & \hrulefill\\
Committee Member & Date\\[0.5in]
\hrulefill & \hrulefill\\
Committee Member & Date
\end{tabular}

\chapter*{Declaration}
\addcontentsline{toc}{chapter}{Declaration}

I declare that this report presents the research and experimental record of the
\systemname\ project as described in the cited repository artifacts. All
claims in the report are bounded by the preregistered protocols, completed
experiments, and immutable evidence indexes. Any institutional declaration
required by the degree-awarding university should replace this provisional
text before submission.

\chapter*{Acknowledgments}
\addcontentsline{toc}{chapter}{Acknowledgments}

The author may insert acknowledgments to supervisors, collaborators,
institutions, funding bodies, and family here. This placeholder is deliberately
non-specific so that no contribution or funding claim is implied.

\chapter*{Abstract}
\addcontentsline{toc}{chapter}{Abstract}

Open-world recognition systems must do more than reject unfamiliar inputs.
They must accumulate rejected evidence, identify persistent candidate groups,
request semantic supervision, integrate confirmed classes, route future
inputs, and preserve auditability and rollback. Research commonly evaluates
these stages independently, but strong stage-level performance does not imply
that their outputs contain the information required by downstream stages.

This thesis investigates stage composition in \systemname, an experimental
framework for explicit geometric class models and transactional model
lifecycles. The realized experiments use frozen 384-dimensional DINOv2
features on bounded CIFAR-10 open-world episodes. The research progresses
through three questions. First, can explicit geometric heads compete with
standard discriminative controls under a frozen representation? Second, can
independently qualified rejection, discovery, confirmation, adaptation, and
routing mechanisms compose into a useful lifecycle? Third, does replacing
stage metrics with post-integration adaptation utility resolve the composition
failure?

The predictive study found that a weighted rank-32 affine head reached 91.73\%
balanced accuracy, remaining 5.03 percentage points behind a 96.77\% RBF SVM
control. The open-world stage study retained a rank-32 class-conditional
Gaussian with 92.08\% known coverage, 86.17\% unknown recall, 0.9561 AUROC,
and 98.67\% review precision. HDBSCAN reached 100\% stage-wise group recall
and review precision, while confirmed rank-16 insertion improved isolated
new-class success by 43.33 points. These stages nevertheless failed to
compose: grouped loops integrated 0 of 3 confirmable classes.

The successor protocol registered adaptation utility---the change in balanced
accuracy on known and discovered classes after integration---as the primary
endpoint at 50 reviewed labels per episode. Density-core review sets covered
only 34.0--39.1\% of the rejected region under a nearest-neighbor support
diagnostic, compared with 99.37--99.78\% for boundary-inclusive sets. Across
nine seed-by-arrival cells, utility selection beat core selection in all nine
by 3.593 points on average, with a paired 95\% bootstrap interval of
[2.794, 4.382]. The result was below the preregistered 5-point practical
threshold, and six cells violated the 2-point remaining-unknown-recall band.
Parent-scoped residual fusion repaired a prior locality defect, preserving
100\% of unaffected predictions, but added at most 0.204 points of utility.

The final outcome is a rigorous negative systems result. Pure discovery cores
are insufficient statistics for adaptation; broader support improves utility;
but representativeness alone does not preserve the remaining unknown space.
The thesis contributes an adaptation-utility endpoint, typed stage-interface
contracts, deterministic episode replay, confirmation-gated transactions,
exact rollback, and artifact-only verification. It concludes that open-world
recognition must be evaluated as a jointly constrained lifecycle rather than
as a sequence of independently optimized benchmarks.

\tableofcontents
\listoftables

\chapter*{List of Abbreviations}
\addcontentsline{toc}{chapter}{List of Abbreviations}

\begin{longtable}{>{\bfseries}p{1.0in}p{4.45in}}
AUROC & Area under the receiver operating characteristic curve\\
BA & Balanced accuracy\\
CSG & Constructive solid geometry\\
DDU & Deep Deterministic Uncertainty\\
EVM & Extreme Value Machine\\
GCD & Generalized category discovery\\
GEODE & Geometric, explicit, open-world development environment used in this work\\
NLL & Negative log likelihood\\
OOD & Out-of-distribution\\
SDF & Signed distance field\\
\end{longtable}

\clearpage
\pagenumbering{arabic}

\chapter{Introduction}

\section{Motivation}

Most classification benchmarks assume a closed world: training and evaluation
share a fixed class inventory, and every input belongs to one of those classes.
Deployed systems rarely enjoy that assumption. New products, species,
documents, faults, behaviors, and domains appear after training. A system that
must always choose a known label can become confidently wrong precisely when
human attention is most valuable.

Open-set recognition addresses part of this problem by allowing rejection.
Open-world recognition extends the objective: a system should identify known
classes, reject unfamiliar inputs, acquire supervision for novel classes, and
incrementally expand its knowledge \cite{bendale2015}. This formulation is
operationally demanding. Rejection creates an unresolved observation, not a
new class. The system must decide what to retain, which rejected examples form
a coherent candidate, which examples deserve limited review, how confirmation
authorizes mutation, how a child model changes acceptance thresholds, and how
the change can be reversed.

Each stage has familiar local metrics. Rejection is summarized by AUROC,
known coverage, and unknown recall. Discovery is summarized by cluster purity,
adjusted Rand index, or category recall. Adaptation is summarized by target
accuracy or forgetting. Routing is summarized by top-$k$ inclusion and compute
reduction. These metrics are useful diagnostics, but they do not answer the
system-level question: did the reviewed evidence create a better model without
damaging what the system still does not know?

\section{Research problem}

The central problem of this thesis is \emph{interface sufficiency}. A producer
stage may optimize a metric while discarding statistics required by its
consumer. A clusterer can produce a pure density core that is easy to label,
while an adapter needs support over the class boundary and low-density modes.
A router can include the correct expert almost always, while hiding exactly
the unknown cases that require exhaustive fallback. A transaction can roll
back exactly while its forward edit changes predictions outside the declared
region.

The thesis therefore treats open-world recognition as a constrained lifecycle
whose interfaces are explicit experimental objects.

\section{Research questions}

The work addresses four research questions:

\begin{description}
\item[RQ1:] Can compact explicit geometric class artifacts approach the
predictive quality of matched discriminative controls under a frozen deep
representation?

\item[RQ2:] Do independently qualified rejection, discovery, human
confirmation, adaptation, and routing stages compose into a review-efficient
open-world loop?

\item[RQ3:] Does selecting reviewed evidence for downstream adaptation utility,
rather than density-core purity, improve post-integration performance at an
equal human-review budget?

\item[RQ4:] Can transactional controls---typed interfaces, confirmation
lineage, edit locality, exact replay, and rollback---make lifecycle failures
auditable without overstating predictive success?
\end{description}

\section{Thesis statement}

The thesis advances the following bounded statement:

\begin{quote}
Stage-wise qualification is insufficient for open-world lifecycle
qualification. At a fixed review budget, representative review selection
improves adaptation utility over density-core selection, but safe lifecycle
advancement additionally requires explicit preservation of known performance,
remaining-unknown recall, lineage, locality, and rollback.
\end{quote}

This statement does not claim closed-set parity, autonomous semantic class
creation, SDF necessity, authoritative learned routing, or successful
end-to-end qualification.

\section{Contributions}

The research makes five contributions:

\begin{enumerate}
\item A controlled comparison separating the predictive value of explicit
geometry from its auditability and editability.

\item A preregistered stage-to-system study showing that strong rejection,
discovery, and isolated adaptation results can fail to compose.

\item A registered adaptation-utility endpoint at a fixed review budget, with
known-class and remaining-unknown safety constraints.

\item A quantified sufficient-statistics diagnosis showing that density cores
are pure but geometrically unrepresentative in the realized episodes.

\item A transactional evaluation harness with typed interfaces, immutable
lineage, deterministic replay, human-confirmation gates, local edit scopes,
and exact rollback.
\end{enumerate}

\section{Organization}

Chapter 2 reviews open-world recognition, discovery, routing, explicit
geometry, and lifecycle safety. Chapter 3 formalizes the problem and endpoint.
Chapter 4 presents the \systemname\ architecture. Chapter 5 records how the
methodology evolved, including failed hypotheses and the tests performed.
Chapters 6 and 7 describe the v7 and v8 experimental protocols. Chapter 8
reports results. Chapter 9 discusses the findings and answers the research
questions. The remaining chapters cover validity, ethics, reproducibility, and
conclusions. Appendices record gates, branch dispositions, and artifact
commands.

\chapter{Background and Related Work}

\section{Open-set and open-world recognition}

Bendale and Boult formalized open-world recognition through explicit rejection
and incremental class addition \cite{bendale2015}. Their Nearest Non-Outlier
method controls open-space risk using metric distance to class means. OpenMax
later calibrated deep activation vectors using extreme-value theory
\cite{bendale2016}. The Extreme Value Machine models Weibull tails over class
margins, yielding incrementally updatable class boundaries \cite{rudd2018}.

Feature-space baselines remain strong when representations are frozen.
Maximum softmax probability is a common baseline \cite{hendrycks2017}.
Mahalanobis scoring uses class-conditional feature distributions
\cite{lee2018}; DDU combines deterministic features with density estimation
\cite{mukhoti2023}; and deep nearest-neighbor distance can provide strong OOD
detection without a parametric head \cite{sun2022}. DUQ and SNGP introduce
distance awareness into discriminative models
\cite{vanamersfoort2020,liu2020}.

These methods play a role similar to the SDF in the original \systemname\
motivation: they define a calibrated acceptance region with an outside signal.
The literature suggests that the mechanism should be treated as substitutable.
This thesis follows that implication by retaining a low-rank Gaussian when it
outperforms geometric support.

\section{Streaming novelty and emerging classes}

The data-stream literature operationalized novelty buffering and class
evolution before modern self-supervised representations. ECSMiner handles
classification and novel-class detection under drift \cite{masud2011}. SAND
and ECHO address scarce labels and evolving concepts
\cite{haque2016sand,haque2016echo}. MINAS maintains micro-cluster summaries
and recurring novelty \cite{faria2016}, while SENCForest targets streaming
emerging classes \cite{mu2017}.

This lineage demonstrates that reject-buffer-group-update loops are not novel
in concept. The open problem examined here is the composition of modern frozen
features, explicit human publication authority, transactional adaptation, and
cross-model support routing under one measured endpoint.

\section{Novel and generalized category discovery}

Generalized category discovery clusters unlabeled data containing known and
novel classes, often over strong self-supervised features \cite{vaze2022}.
ORCA jointly learns seen and novel categories in open-world semi-supervised
learning \cite{cao2022}. Grow and Merge and incremental GCD extend discovery
across episodes \cite{zhang2022,zhao2023}.

Clustering mechanisms include DBSCAN \cite{ester1996}, HDBSCAN
\cite{campello2013}, and FINCH \cite{sarfraz2019}. Purity and category recovery
are natural discovery metrics, but they do not measure whether a small reviewed
subset supports a downstream class model. The present study isolates this
difference.

\section{Continual learning and modular routing}

Continual learning studies acquisition under stability-plasticity constraints.
The present work uses a narrower transactional setting: the representation is
frozen, semantic mutation requires confirmation, and every child has an exact
parent.

Expert Gate routes examples among task experts using reconstruction error
\cite{aljundi2017}. DEMix and Branch-Train-Merge provide modern modular
precedents \cite{gururangan2021,li2022}. In \systemname, fingerprints establish
structural compatibility, while empirical profiles estimate support. A
profile may only become authoritative if it agrees with exhaustive inference
and fails safely on unknown inputs.

\section{Explicit geometry and low-rank class models}

Signed distance fields and constructive solid geometry provide editable
spatial representations \cite{requicha1980,park2019}. Their attraction in
feature space is governance: an operator could inspect components, remove a
lobe, merge regions, or roll back an edit. Predictive competitiveness,
however, is a separate empirical question.

The retained affine and Gaussian components use a center, orthonormal tangent
basis, tangent variances, and residual variance. This construction relates to
probabilistic PCA and mixtures of factor analyzers
\cite{tipping1999,ghahramani1996}. Proper likelihood includes both normalized
distance and volume terms; raw geometric support does not automatically yield
cross-class comparability.

\section{Research gap}

Prior work validates each stage individually. What is uncommon is an
end-to-end study in which:

\begin{itemize}
\item review cost is fixed;
\item stage interfaces and sufficient statistics are registered;
\item semantic mutation requires linked confirmation;
\item child thresholds and routes carry lineage;
\item exact rollback and locality are independently tested; and
\item the primary endpoint is post-integration utility.
\end{itemize}

This thesis addresses that systems-level gap. Its novelty claim concerns the
controlled composition study and audit boundary, not the existence of the
discover-label-update loop or any individual mechanism.

\chapter{Problem Formulation}

\section{Open-world episode}

Let $f(x)\in\mathbb{R}^d$ be a frozen representation. At the beginning of
episode $e$, parent model $M_e$ recognizes class set $\mathcal{K}_e$. The
stream contains known inputs and an arriving class $c_e\notin\mathcal{K}_e$;
a separate set $\mathcal{R}_e$ remains unknown after the episode.

The rejector maps an input to a known class or rejection. Rejected records enter
a bounded buffer. A discovery policy forms a persistent review candidate.
A review selector chooses at most $B$ examples. An authorized confirmation
assigns semantic meaning, after which an adapter may publish child
$M_{e+1}$. A router chooses compatible bundles or falls back to exhaustive
evaluation.

\section{Semantic states}

Three states must remain distinct:

\begin{description}
\item[Rejected:] an acceptance score crosses a threshold.
\item[Grouped:] rejected examples exhibit persistent structure.
\item[Confirmed:] an authorized review assigns semantic meaning.
\end{description}

Only confirmation authorizes a semantic class or model mutation. Group IDs are
persistent review objects, not provisional semantic labels.

\section{Adaptation utility}

The v8 primary endpoint is episode utility:

\begin{equation}
\utility_e =
\ba_{\mathcal{K}_e\cup\{c_e\}}(M_{e+1})
-
\ba_{\mathcal{K}_e\cup\{c_e\}}(M_e).
\label{eq:utility}
\end{equation}

The parent receives zero credit on the unavailable arriving class. This avoids
the misleading comparison in which a no-update parent is evaluated only on
classes it already knows.

The cumulative endpoint after $N$ episodes is

\begin{equation}
\utility_{1:N} =
\ba_{\mathcal{K}_1\cup\{c_1,\ldots,c_N\}}(M_{N+1})
-
\ba_{\mathcal{K}_1\cup\{c_1,\ldots,c_N\}}(M_1).
\end{equation}

\section{Constraints}

Utility is evaluated subject to:

\begin{align}
\text{reviews}_e &\leq B,\\
\ba_{\mathcal{K}_e}(M_e)-\ba_{\mathcal{K}_e}(M_{e+1})
&\leq \epsilon,\\
\operatorname{Recall}_{\mathcal{R}_e}(M_e)-
\operatorname{Recall}_{\mathcal{R}_e}(M_{e+1})
&\leq \delta.
\end{align}

The registered values are $B=50$, $\epsilon=0.01$, and $\delta=0.02$.
Transactions must also satisfy confirmation, graph, lineage, replay, rollback,
and routing-fallback invariants.

\section{Interface sufficiency}

For adjacent stages $P\rightarrow C$, let $S_C$ be the sufficient statistics
needed by consumer $C$ and $E_P$ be the evidence emitted by producer $P$. An
interface mismatch exists when

\begin{equation}
S_C \nsubseteq E_P.
\end{equation}

The central candidate mismatch is clusterer-to-review: the producer emits dense
core members and persistence, while the adapter needs support over modes,
covariance directions, and parent-boundary regions. The experiments compare
review policies that alter $E_P$ at fixed cost.

\section{Registered decisions}

The pivotal utility-selected policy advances only if its mean gain over core
selection is at least 5 points, the paired 95\% bootstrap lower bound is above
zero, at least 7/9 cells are positive, and all safety and transaction
constraints pass. The 5-point operand is a practical effect threshold. It is
not replaced by a positive confidence interval after results are known.

\chapter{The \systemname\ System}

\section{Design objectives}

\systemname\ separates five concerns:

\begin{enumerate}
\item frozen representation;
\item class-conditional support geometry;
\item calibrated decision or rejection;
\item versioned lifecycle operations; and
\item evidence and reproducibility.
\end{enumerate}

This separation allows a head to fail predictive parity while its transaction
machinery remains useful. It also permits a Gaussian or kNN mechanism to
replace an SDF without weakening a systems claim.

\section{Explicit class components}

For a component with center $\mu$, orthonormal tangent basis $V\in
\mathbb{R}^{d\times r}$, tangent variance vector $\lambda$, and isotropic
residual variance $\sigma^2$, define

\begin{equation}
q(x)=
\sum_{j=1}^{r}\frac{(v_j^\top(x-\mu))^2}{\lambda_j}
+
\frac{\lVert(I-VV^\top)(x-\mu)\rVert_2^2}{\sigma^2}.
\end{equation}

The corresponding Gaussian log likelihood is

\begin{equation}
\log p(x\mid k)=-\frac{1}{2}
\left[
q(x)+\sum_{j=1}^{r}\log\lambda_j+
(d-r)\log\sigma^2+d\log(2\pi)
\right].
\end{equation}

The volume terms are essential for cross-class comparison. A raw field that
captures coverage can still be poorly calibrated for discrimination.

\section{Immutable model bundles}

A bundle records the representation identity, class order, components,
readout, thresholds, router profile, parent hash, and schema version.
Publication creates an immutable child rather than mutating the parent in
place. Rollback changes the current pointer to the exact parent while retaining
the full append-only history.

\section{Review and confirmation}

Rejected observations carry stable record IDs and representation lineage.
Discovery produces persistent review IDs. Confirmation records reference those
reviews and contain the authorized outcome. The adapter rejects missing,
stale, or mismatched confirmation. This design prevents a coherent unlabeled
cluster from silently becoming a semantic class.

\section{Typed interfaces}

The v8 protocol registers five interfaces:

\begin{enumerate}
\item rejector to buffer;
\item buffer to clusterer;
\item clusterer to review selector;
\item review to adapter; and
\item adapter to router.
\end{enumerate}

Seven versioned schemas encode endpoint evidence, episode replay, interface
audits, threshold transfer, review selection, residual scope, and integrated
routing. Six negative cases verify fail-closed behavior for leakage, stale
class order, threshold mismatch, budget overflow, missing confirmation, and
rollback-parent drift.

\section{Research progression}

The complete program is cumulative. Early studies tested primitive fitting,
constructive operations, calibrated readouts, precision policies, scaling, and
transactions. The thesis focuses on the progression most relevant to its
questions:

\begin{description}
\item[v6/v6.1:] predictive competitiveness of explicit heads and lifecycle
locality;
\item[v7:] stage-wise open-world qualification and integrated composition;
\item[v8:] adaptation utility, interface statistics, equal-budget review
selection, and parent-scoped residuals.
\end{description}

Every milestone uses immutable parents and registered advancement gates.
Failed gates close downstream branches rather than being silently redefined.

\chapter{Evolution of the Research Methodology}

\section{Why the methodological history matters}

The final adaptation-utility protocol was not the initial methodology. It
emerged after a sequence of increasingly restrictive experiments exposed
different failure modes. Recording only the final protocol would hide why its
endpoint, controls, and stopping rules exist. This chapter therefore separates
the \emph{research process} from the final experimental methods.

The progression can be summarized as:

\begin{equation}
\begin{aligned}
\text{implementation correctness}
&\rightarrow \text{geometric expressiveness}\\
&\rightarrow \text{predictive competitiveness}\\
&\rightarrow \text{stage qualification}\\
&\rightarrow \text{lifecycle utility}.
\end{aligned}
\end{equation}

Each transition occurred because the previous level proved necessary but
insufficient. The methodology did not advance by treating every experiment as
a success. Null, negative, blocked, and infeasible outcomes determined which
questions were legitimate to ask next.

\section{Phase I: implementation correctness and semantic tests}

The recorded methodology begins with the repaired Tier-6 temporal baseline and
M0 protocol lock, followed by M1 semantic correctness. The earliest methodology
focused on whether the implemented geometric
operations meant what their formulas and APIs claimed. These tests preceded
performance comparisons because a predictive result is uninterpretable when
the primitive semantics are wrong.

The test families included:

\begin{itemize}
\item seeded-RNG ownership tests, which removed dependence on ambient NumPy
state;
\item normalized-softmin conservation tests, which detected an exact
$\log(4)/2$ field shift under pruning;
\item finite-difference gradient tests for hard constructive operations;
\item scalar-versus-vectorized optimizer parity;
\item CPU-versus-OpenCL numerical parity;
\item serialization and constructor smoke tests; and
\item numerical audits of approximate ellipsoid distance.
\end{itemize}

The ellipsoid audit established an important claim boundary. The implemented
radial correction is exact for spheres but can deviate substantially for
eccentric ellipsoids. Terminology was therefore narrowed from an exact metric
distance claim to an implemented first-order correction. This pattern---test a
semantic claim, then restrict the claim when the test fails---became a core
methodological rule.

\section{Phase II: M0--M4 baseline construction and causal ablation}

M0 froze experiment configurations,
dataset manifests, array fingerprints, seeds, split policies, resource
accounting, and deterministic smoke artifacts. The purpose was to make later
comparisons traceable to identical inputs.

M2 then decomposed geometric fitting from decision readout. Geometry and
calibration used disjoint data. Five readouts were compared on the same fitted
artifact. Temperature scaling left accuracy unchanged at 84.69\% while
reducing NLL from 1.48 to 0.50 and ECE from 0.60 to 0.03; multinomial readout
added 0.19 accuracy points and reduced NLL to 0.49. This established that
field quality and probability calibration had to be measured separately.

M3 introduced eight matched classical baselines over identical transformed
features and split hashes. Over five seeds, logistic regression reached
84.02\%, RBF SVM 84.00\%, and the GEODE multinomial readout 83.11\%.
The methodology thereafter required same-representation controls rather than
comparing against unrelated published numbers.

Early causal ablations examined additive geometry, held-out subtractive
excision, and active repair. Subtractive constructive geometry showed null or
slightly negative aggregate effects, and active repair generated no accepted
candidates. The methodology therefore stopped treating more expressive CSG as
automatically beneficial. Subtraction remained opt-in rather than becoming a
default source of unmeasured complexity.

The testing strategy expanded from unit semantics to:

\begin{itemize}
\item matched ablations with identical representations and budgets;
\item seed-level and pooled paired intervals;
\item corruption and leakage fixtures;
\item runtime, memory, and serialized-state accounting; and
\item append-only evidence manifests.
\end{itemize}

\section{Phase III: M5--M7 stress, sequence, and uncertainty studies}

M5 asked whether alternative high-dimensional fitters improved robustness
under fixed budgets. Seven synthetic fitters passed through a synthetic gate,
a five-seed Tier-4 screen, and a locked Tier-5 confirmation. A shrinkage
candidate gained 0.75 points on the selection screen but lost 1.56 points on
confirmation, used 4.4 times as many primitives, and took 3.2 times longer.
The current fitter was retained. This result motivated confirmation stages and
resource co-primary metrics.

M6 tested whether explicit geometric state could represent temporal context.
Exact, reservoir, ensemble, and hybrid causal representations were compared
through periodic synthetic gates, train-only tuning, and a locked
50,000/10,000 WikiText episode. GEODE reached 30.36\%, below a linear model at
34.64\%, a matched 5-gram at 44.50\%, and a practical 5-gram at 47.61\%.
The result established representational flexibility, not competitive temporal
learning.

M7 expanded evaluation beyond top-1 accuracy to controlled shift, selective
prediction, calibration, conformal diagnostics, and real CIFAR-100/SVHN OOD
families. Maximum-probability AUROC reached 0.786 and 0.831 and exceeded the
tested Mahalanobis and kNN scores, but FPR95 remained 0.579 and 0.431. The claim
was restricted to relative ranking under the frozen fit.

\section{Phase IV: M8--M10 editability, scaling, and open-set inference}

M8 combined corruption tests, seed-stability tests, explicit insertion and
deletion, rollback, and online updates. Broad robustness was rejected, while
predictions remained 93.0\% seed-agreeing despite primitive drift. Behavioral
edit gates showed that insertion, deletion, and rollback could be local and
reversible in the tested setting.

M9 measured editability rather than assuming that compact artifacts scaled.
One-factor sweeps varied class count, dimension, and primitives per class over
five repeats with operation counters. Locality gates passed and insertion
stayed below 0.18 ms, but exhaustive inference and dense 128-dimensional
snapshots exposed bounded scaling costs.

M10 introduced support profiles, abstention, and leave-class-out selection in
a locked three-episode toy matrix. No score, threshold, and geometry
combination jointly reached 90\% known coverage and 50\% unknown recall across
all episodes. The experiment showed that calibrated closed-set scores did not
automatically produce reliable open-set inference.

\section{Phase V: M11 discovery methodology}

M11 was not one experiment but a long methodological sequence on accumulated
unsupported observations. It began with bounded evidence, automatic
clustering, stable review IDs, and confirmation-gated adaptation. Initial
studies exposed category merging and fragmentation, after which the sequence
tested:

\begin{itemize}
\item DBSCAN flag budgets and radius selection;
\item HDBSCAN under variable-density hypotheses;
\item PCA versus normalized raw MobileNet representations;
\item corrected evaluation on the final accumulated partition rather than the
first review snapshot;
\item larger accumulation windows;
\item FINCH without a supplied group count;
\item DenStream-style fading micro-clusters;
\item silhouette-selected $k$-means;
\item transductive joint PCA;
\item delayed must-link and cannot-link metric learning;
\item COBRAS-style split and merge feedback;
\item answer-blind active pair queries;
\item missing, flipped, and contradictory-answer stress tests;
\item independent confirmation of pairwise answers;
\item 450-observation lower-tail reliability studies at each error rate; and
\item a ResNet-18 representation control.
\end{itemize}

The evaluation boundary itself evolved. Early group metrics used first-review
snapshots, which did not match the operational goal. M11.26 separated online
review delay from autonomous partition quality and evaluated the complete
flagged buffer, including noise. Under the corrected protocol, the retained
normalized-MobileNet HDBSCAN policy reached 61.1\% distinct-group recall,
79.8\% purity, and full two-group recovery in 4/9 cells, but noise-inclusive
adjusted Rand index was only 0.059.

Feedback studies showed positive mean improvements but weak lower-tail safety.
At 12.5\% respondent error, independently confirmed constraints still accepted
harmful constraints in 10.44\% of observations; fifth-percentile ARI change was
negative. Constraints therefore remained analysis-only and unpersisted.
M11.40 closed the sequence by retaining only the automatic review-grouping
policy and prohibiting autonomous mutation.

\section{Phase VI: M12 routing and scaling methodology}

M12 began with one-factor synthetic sweeps over classes, dimensions, primitive
counts, and batch size. Exhaustive inference exhibited approximately linear
latency growth in class and primitive count. Subsequent iterations tested:

\begin{itemize}
\item conservative support-sphere lower bounds with exact early stopping;
\item grouped active sample-class evaluation;
\item incumbent initialization;
\item nearest-neighbor indexes over support centroids;
\item greedy primitive removal;
\item disjoint confirmation streams for tentative pruning; and
\item a complete exit audit with counters, agreement, fallback, and latency.
\end{itemize}

Some candidates preserved exact predictions, but they either evaluated too
much of the exhaustive class set, failed normalized-score requirements, or did
not improve tail latency. The methodology therefore rejected operation-count
improvement as sufficient evidence for production routing.

\section{Phase VII: M13--M15 primitive and field semantics}

M13 reopened primitive-family comparisons under matched full-covariance,
diagonal-covariance, and trace-matched spherical fits, including
primitive-aware OpenCL paths and the full additive/excision/repair matrix.
Spherical additive geometry reached 83.80\%, 0.463 points above diagonal, but
subtraction again added no benefit and remained disabled.

M14 reinterpreted frozen covariance primitives as hierarchical Gaussian
mixtures. Proper probabilistic fields substantially improved NLL and ECE but
not accuracy; for example, sphere multinomial NLL improved from 0.4847 to
0.4772 while accuracy moved by $-0.15$ points. This reinforced the separation
between calibration and classification.

M15 tested auditable fusion of geometric and probabilistic field scores.
Global and per-class covariance temperatures and hybrid readouts were fit under
a matched five-seed protocol. Hybrid NLL improved over either field-only input
within the accuracy tolerance, but a feature-space logistic control remained
better. These findings motivated the later use of strong substitutable controls
instead of making explicit-field identity the endpoint.

\section{Phase VIII: E0--E11 systems qualification before research claims}

The project next added resumability, content-addressed bundles, compatibility
validation, activation, and rollback. This phase did not claim predictive
superiority. Its question was whether experiments and model changes could
survive interruption and be reproduced exactly.

Failure-injection tests covered:

\begin{itemize}
\item class-fitting interruption and partial-class replacement;
\item optimizer, momentum, scheduler, scaler, and RNG restoration;
\item per-epoch cursor and distributed-sampler restoration;
\item incompatible worker-topology rejection;
\item corrupted bundle and fingerprint mismatch rejection;
\item failed canaries and coordinator loss; and
\item deterministic JSON/CSV history export.
\end{itemize}

These tests established the parent-child transaction model later used for
adaptation. They also established a distinction that persists throughout the
thesis: a mechanism may qualify operationally without qualifying
predictively.

\section{Phase IX: v5 metric and representation falsification}

The v5 methodology introduced versioned stages, deterministic factorial
registries, lineage guards, paired statistics, Pareto comparisons, and
canonical result artifacts. Ten representation-by-head cells replayed
byte-identically, while a deliberate split mismatch failed closed.

Six precision families---spherical, diagonal, full, low-rank, shared, and
policy-selected variants---were then tested. The adaptive metric-support policy
lost 0.323 balanced-accuracy points and reduced neither parameter bytes nor
deterministic fitting work. Both predictive and efficiency gates failed.

This negative result changed the research strategy. Rather than adding more
unregistered metric flexibility, the next program locked a strong
representation and matched controls, then asked which geometric choices could
close a measured predictive gap.

\section{Phase X: v6 predictive methodology}

\subsection{Locked parent and cheap falsification}

M27 froze parent evidence, baseline predictions, teacher lineage, primitive
budgets, and three evaluation seeds. The methodology used an S1 cheap
falsification stage before the more expensive S2 confirmation stage.

Boundary distillation into component-matched spheres was the first candidate.
It reached only 79.80\% development accuracy, 6.30 points below the retained
explicit model, and stopped at S1 despite exact replay. This rejected the
hypothesis that imitating the RBF boundary with the existing primitive family
was sufficient.

\subsection{Subspaces and directional geometry}

Ranked affine-subspace primitives were introduced next. Rank-32 radial scoring
reached 89.00\%, improving the prior explicit head by 2.90 points and passing
the cheap gate. Spherical caps then tested whether directional geometry
explained additional error. Cosine caps beat matched Euclidean spheres by
3.97 points over seeds 11, 23, and 37, with a positive seed-level interval.

These results supported two intermediate conclusions: frozen deep features
contain useful low-rank and directional structure, and primitive choice matters.
They did not establish parity with the discriminative control.

\subsection{Identifiable factorial and parity stopping rule}

M31 crossed objective, primitive, score, and budget factors in an identifiable
fractional factorial. The direct rank-32 radial family reached 90.77\%, but
remained 6.00 points behind the RBF SVM. Because the registered parity operand
failed, downstream parity-rescue branches closed.

The v6.1 amendment allowed only bounded, mechanistically motivated tests:
tangent caps, constrained nonnegative weights, a component-budget diagnostic,
an independent locality frontier, and a separate few-shot feasibility track.
Weighted affine readout improved the explicit model to 91.73\%, but the final
5.03-point RBF deficit still exceeded the 2-point gate. The predictive program
ended with Outcome D rather than extending the search indefinitely.

\section{Phase XI: lifecycle claims separated from accuracy claims}

The A3 edit suite tested 24 explicit transactions. Every transaction rolled
back exactly, but unaffected-prediction preservation fell to 84.2\% against a
99.9\% locality contract. This result showed that exact rollback and local
forward behavior are different properties.

The Flowers few-shot track separately tested feasibility. Rank-3 fitting was
possible with five examples per class and replayed exactly; rank 32 required
34 examples per class and was declared infeasible rather than estimated from
insufficient support. This introduced a further methodological category:
\emph{infeasible} is distinct from \emph{failed}.

\section{Phase XII: v7 stage-wise open-world methodology}

With predictive parity closed, the research question changed. The claim was no
longer that explicit geometry matches an RBF classifier. It became whether an
auditable open-world loop could qualify stage by stage and compose.

M38 preregistered the claim boundary, immutable parents, schedule families, and
semantic authority rules. M39 compared six matched rejection heads and retained
the low-rank Gaussian, demonstrating that the system claim did not depend on
retaining an SDF. M40 compared five grouping policies. M41 tested confirmation-
gated updates. M42 tested empirical routing against exhaustive authority.
M43 then admitted only stage-passing mechanisms into the integrated factorial.

This sequential retention design reduced researcher degrees of freedom:

\begin{enumerate}
\item a mechanism had to pass its local gate before composition;
\item failed routing remained shadow-only;
\item semantic publication always required confirmation;
\item integrated failure blocked untouched confirmation; and
\item stage evidence could not be relabeled as lifecycle evidence.
\end{enumerate}

The final composition failure changed the methodology again. Purity,
unknown-recall, and isolated target-gain gates had all appeared successful, but
the composed loop integrated no grouped class.

\section{Phase XIII: v8 endpoint redesign}

The v8 program was registered as a new claim rather than a rescue of v6.1 or
v7. Its guiding change was to replace per-stage advancement as the primary
criterion with post-integration adaptation utility at a fixed review budget.

M45 formalized interface contracts and six fail-closed mismatch tests. M46
tested threshold transfer and audited sufficient statistics. M47 compared six
review policies over nine paired cells. M49 independently tested whether
parent-scoped fusion repaired locality.

The pivotal design included explicit kill switches. M48 learned selection and
M50 untouched lifecycle qualification could run only if M47 passed practical,
statistical, safety, and transactional gates. M47 produced a consistent
positive effect but failed the practical and safety operands. The downstream
branches therefore remained closed.

\section{Phase XIV: the joint-budget scaling frontier}

The final programme phase moved from the stage-wise open-world setting to a
controlled measurement of the frozen patch-dictionary classifier's scaling
behaviour. The system under study is a deterministic encoder (whitened patch
dictionary plus pooling) feeding a closed-form ridge head --- no learned
backbone, no iterative optimisation. Two resources are available: dictionary
size (the width of the head's feature space) and training data. All results
are registered before measurement, gated by kill switches, anchored by exact
reproduction of prior sealed numbers, and compared at matched per-image
multiply-accumulate cost on a fixed 34{,}500-row DomainNet test set.

The central finding rescinds an earlier negative. The single-stage accuracy
``ceiling'' was a slice artifact: measured on the two-dimensional
(dictionary-size $\times$ data) surface, the family is super-additive ---
scaling both axes jointly beats the sum of scaling either axis alone
(measured excess +0.0097/+0.0101 over a +0.005 margin), and 6{,}144-atom
dictionaries reach 0.2249 at full data versus 0.2153 at 3{,}072 atoms. The
compute is data-elastic: wider heads consume data more productively, with
per-atom learning-curve steepness rising 0.094, 0.115, 0.134 across the atom
ladder. The fitted learning-rate exponent is width-independent (gamma about
$-0.012$), so data-elasticity is an absolute-gain effect, not a faster rate.

Two structural limits were then measured. First, the atoms axis saturates:
extending the dictionary past about 6{,}144 atoms at full data yields a flat,
slowly declining ridge (0.2246 to 0.2205 from 6{,}144 to 16{,}384 atoms), so
the remaining lever is data, not dictionary size. Second, the gap to a frozen
DINOv2 trunk (0.22 versus 0.54 at the large end) is structural: the code
space is dominated by a fixed eight-dimensional structure (effective rank
7.8, flat across all dictionary sizes), and the correct class loses the
argmax for roughly 78\% of test samples --- accuracy lives in a thin
positive-margin tail. A spectral margin certificate was registered and
failed (it cannot predict the argmax crossing) and is retained as an
explanatory diagnostic only.

The v19 phase closed four routes (closed-form depth, accuracy-preserving
binary codes, spectral certificates, and dictionary growth past 6{,}144
atoms) and opened one: the data axis, which is the steep one and the only
remaining joint-budget lever. A follow-on engineering track (v20) measured
the system-architecture implications: contract-gated routing of programmatic
primitives rejects out-of-contract inputs with zero learned forward passes
while preserving in-contract accuracy, and a fully programmatic variable-
order memory (the ``how far back'' dial) shows a measured perplexity optimum
as a next-token model on a constrained DSL, with footprint growing sharply
past the optimum window. Details and reproduction steps are in the sealed
evidence (\texttt{RESEARCH\_REPORT\_v19}, \texttt{ENGINEERING\_PLAN\_v20}).

\section{Experiment and test ledger}

Table~\ref{tab:evolution-ledger} records the principal hypothesis-driven
experiments that led to the final thesis. ``Tests'' denotes focused regression
tests attached to that milestone where recorded; factorial cells and seeds are
reported separately because they are experimental observations, not software
unit tests.

\begin{longtable}{
>{\raggedright\arraybackslash}p{0.6in}
>{\raggedright\arraybackslash}p{1.45in}
>{\raggedright\arraybackslash}p{0.9in}
>{\raggedright\arraybackslash}p{1.95in}}
\caption{Evolution of the principal methodology and tests.}
\label{tab:evolution-ledger}\\
\toprule
\textbf{Stage} & \textbf{Question or intervention} & \textbf{Design} &
\textbf{Decision}\\
\midrule
\endfirsthead
\toprule
\textbf{Stage} & \textbf{Question or intervention} & \textbf{Design} &
\textbf{Decision}\\
\midrule
\endhead
Tier 6 & Repair temporal state, calibration, and GPU semantics & WikiText real
run and pipeline benchmark & Established 23.91\% isolated and 23.35\% pipeline
baselines; no broad sequence claim.\\
M0 & Freeze protocols, hashes, manifests, and resource counts & Canonical
configs and smoke artifacts & Reproducible evidence base; no accuracy claim.\\
M1 & Are primitive semantics and implementations correct? & 28 then 34 fast
tests & RNG, pruning, CSG gradients, and CPU/GPU parity repaired.\\
M2 & How much error belongs to calibration versus geometry? & 5 readouts on one
split & Temperature and multinomial readouts sharply improved NLL/ECE.\\
M3 & How does GEODE compare on identical features? & 8 baselines $\times$ 5
seeds & GEODE trailed logistic by 0.91 points.\\
M4 & Do subtractive CSG and active repair help causally? & 3 ablations $\times$
5 seeds & Intervals included zero; active repair added no candidates.\\
M5 & Which high-dimensional fitter survives confirmation? & 7 fitters; Tier-4
screen and Tier-5 lock & Screen winner reversed on confirmation; current fitter
retained.\\
M6 & Can explicit state model temporal context competitively? & Synthetic gates
and 50k/10k WikiText & Flexibility shown; linear and n-gram controls won.\\
M7 & Are confidence and OOD scores reliable? & Controlled shift plus
CIFAR-100/SVHN & Relative score ranking established; FPR95 limited claim.\\
M8 & Are models robust, stable, and behaviorally editable? & 5-seed corruption,
stability, edit tests & Broad robustness rejected; tested edits reversible.\\
M9 & How does editability scale? & Class, dimension, and primitive sweeps; 5 repeats &
Locality passed; exhaustive and snapshot costs exposed.\\
M10 & Can calibrated scores support open-set inference? & 3 locked episodes &
No cell passed known and unknown gates jointly.\\
M11 & Can accumulated rejects form safe persistent groups? & 40 iterative
studies and closure audit & Partial group recovery; semantics and mutation
remained human-gated.\\
M12 & Can exact routing become sublinear and faster? & 8 scaling, bound,
indexing, pruning, and audit iterations & Exact candidates did not satisfy
score, work, and tail-latency requirements jointly.\\
M13 & Which primitive family and CSG policy helps? & 5-seed CPU screen and
15-cell GPU matrix & Sphere additive retained; subtraction again null.\\
M14 & Do probabilistic field semantics improve calibration? & Matched
CPU and OpenCL comparisons & NLL and ECE improved; accuracy did not.\\
M15 & Does geometric/probabilistic fusion help? & 5-seed matched readouts;
179-test gate & Hybrid NLL improved; feature logistic remained better.\\
E0--E2 & Can interrupted research resume exactly? & 19 focused tests growing
to a 187-test gate & Feature, class, optimizer, sampler, RNG, and history state
recovered exactly.\\
E3--E4 & Can bundles and a real CIFAR lifecycle qualify? & 192 then 196 tests &
Bundle lifecycle passed; CIFAR non-inferiority passed with weak near-OOD FPR95.\\
E7 & Can distributed DomainNet execution qualify? & 3 logical nodes; 192-image
episode & Local recovery passed; physical multi-host gate remained blocked.\\
E10--E11 & Can serving recover and public artifacts reproduce? & Fault
rehearsal and 30 frozen inputs & 1-second RTO/zero-request RPO passed; five
outputs reproduced byte-identically.\\
M16 & Can the protocol replay and reject split drift? & 10 cells; 15 tests &
Byte-identical replay; deliberate split mismatch rejected.\\
M18 & Do adaptive precision families improve accuracy or cost? & 6 families;
13 tests & Lost 0.323 points; no resource reduction; stopped.\\
M27 & Can v6 begin from immutable matched parents? & 3 heads; 254-test full
gate & Parents, predictions, teacher, and budgets locked.\\
M28 & Can RBF boundaries be distilled into matched spheres? & S1; 4 tests &
79.80\%; 6.30 points below parent; stopped.\\
M29 & Do low-rank affine subspaces help? & Rank sweep; 6 tests & Rank-32 radial
gained 2.90 points; advanced.\\
M30 & Does directional geometry matter? & 3 seeds; 6 tests & Cosine exceeded
Euclidean by 3.97 points; mechanism retained.\\
M31 & Which factor family approaches parity? & 9 cells $\times$ 3 seeds &
90.77\%, 6.00 points behind RBF; parity branches closed.\\
A1-T & Do tangent-space caps rescue the model? & S1; 10 tests & Worse than
affine control; stopped.\\
A1-W & Does constrained component weighting help? & S1; 10 tests & Added 2.60
points and reduced NLL; advanced.\\
A1-B & Is more component capacity monotonic? & 8 tests & Peaked then declined;
classified as unstable.\\
A2 & Does weighting close parity on all seeds? & 3 seeds; 7 tests & 91.73\%,
5.03 points behind RBF; Outcome D.\\
A3 & Are edits local and reversible? & 24 transactions; 7 tests & Rollback
exact; locality only 84.2\%.\\
A4 & Is five-shot fitting feasible? & 2 rank tiers; 4 tests & Rank 3 feasible;
rank 32 blocked by sample count.\\
v6.1 final & Can conclusions replay without source data? & 7 indexes, 65
artifacts; 5 tests & 11 conclusions reproduced byte-identically.\\
M38 & Are v7 claims, parents, and schedules frozen? & 3 schedule families; 10
tests & Protocol and prior-art boundary locked.\\
M39 & Which rejection head passes matched gates? & 6 heads $\times$ 3 seeds &
Only rank-32 Gaussian retained.\\
M40 & Which grouping policies create stable review objects? & 5 policies; 15
cells & HDBSCAN primary; FINCH control.\\
M41 & Which confirmed updates are safe? & 24 operation cells & New-class
rank-16 insertion retained.\\
M42 & Can empirical routing become authoritative? & 4 profiles; 12 cells & No;
routing remained shadow-only.\\
M43 & Do retained stages compose? & 36 factorial cells & Grouped arms integrated
0/3 classes; Outcome C.\\
v7 final & Can Outcome C replay artifact-only? & 6 indexes, 16 artifacts &
8 conclusions reproduced twice.\\
M45 & Can interfaces fail closed? & 7 schemas; 6 mismatch cases; 11 tests & All
registered mismatches rejected.\\
M46 & Do thresholds transfer and cores represent support? & 4 rules; 3 seeds;
5 tests & Anchor quantile retained; core support defect identified.\\
M47 & Does utility selection beat core at equal budget? & 6 policies; 54 arm
cells; 3 tests & Positive 9/9, but practical and safety gates failed.\\
M49 & Can parent-scoped residuals repair locality and utility? & 2 attempts; 2
tests & 100\% locality; no residual retained.\\
v8 final & Can Outcome D replay artifact-only? & 4 indexes, 11 artifacts; 1
test & 6 conclusions reproduced twice without data access.\\
\bottomrule
\end{longtable}

\section{Test taxonomy}

The test count alone does not characterize the evidence. The final 374-test
repository gate combines several distinct test classes:

\begin{description}
\item[Semantic unit tests] verify equations, normalization, gradients, shapes,
and deterministic RNG behavior.

\item[Property and invariant tests] verify conservation, orthonormality,
monotonic fail-closed rules, graph consistency, and immutable class order.

\item[Parity tests] compare scalar and vectorized implementations, CPU and GPU
paths, and pre- and post-serialization predictions.

\item[Failure-injection tests] interrupt fitting, corrupt bundles, alter
fingerprints, remove confirmation, stale thresholds, exceed budgets, or lose
workers and coordinators.

\item[Leakage tests] reject split mismatch, partition overlap, use of
validation examples as adaptation support, and final-label access.

\item[Lifecycle tests] exercise publication, child lineage, pointer changes,
rollback, stale-profile fallback, and audit history.

\item[Statistical tests] implement paired comparisons, bootstrap intervals,
seed consistency, and practical effect gates.

\item[Artifact-replay tests] recompute frozen conclusions from indexed JSON
without loading training data or opening sealed labels.
\end{description}

Software tests establish implementation and evidence integrity; they do not
replace experimental replication. Conversely, repeated experimental cells do
not prove that transaction and leakage invariants are implemented correctly.
The methodology deliberately requires both.

\section{How negative results shaped the final design}

The main methodological changes can be traced directly to failures:

\begin{longtable}{
>{\raggedright\arraybackslash}p{1.85in}
>{\raggedright\arraybackslash}p{3.35in}}
\toprule
\textbf{Observed failure} & \textbf{Resulting methodological change}\\
\midrule
\endhead
CSG subtraction was null or negative & Keep subtraction opt-in and require
causal matched ablations.\\
Adaptive metric support reduced neither error nor cost & Stop adding metric
flexibility without a registered utility hypothesis.\\
Explicit heads remained behind RBF & Treat acceptance heads as substitutable;
separate editability from predictive superiority.\\
Rollback passed while locality failed & Register locality independently and
freeze affected regions from the parent.\\
Strong stages failed to compose & Replace stage metrics with adaptation
utility as the primary endpoint.\\
Pure cores failed adaptation & Audit producer-consumer sufficient statistics
and compare representative review policies.\\
Utility gain harmed unknown recall & Make remaining-unknown recall a binding
lifecycle constraint.\\
Positive interval missed practical threshold & Preserve Outcome D and block
post-hoc learned selection and untouched confirmation.\\
\bottomrule
\end{longtable}

\section{Methodological maturity}

The methodology matured from exploratory implementation to preregistered
falsification. Early experiments discovered semantics and established
baselines. Later programs froze parents, controls, budgets, seeds, effect-size
thresholds, safety constraints, and branch dependencies before evaluation.
This evolution reduces flexibility but strengthens the interpretation of a
negative result.

The final thesis is therefore not a report of one successful algorithm. It is
a documented sequence of hypotheses in which each retained conclusion is tied
to a testable gate, and each stopped branch remains part of the scientific
record.

\chapter{Stage-Wise Open-World Study}

\section{Data and representation}

The realized experiment uses immutable 384-dimensional CLS-token features from
DINOv2-small \cite{oquab2024} on CIFAR-10. Each frozen seed contains 10,000
training and 1,000 development examples. Seeds are 11, 23, and 37. In the v7
leave-two-class-out protocol, classes 0--7 are known, class 8 is confirmable,
and class 9 remains unknown.

Broader CIFAR-100 and DomainNet tracks described in earlier planning were not
executed in this chain and are not presented as evidence. Final confirmation
labels remained sealed after the integrated gate failed.

\section{Acceptance-head comparison}

Six heads share the same representation and fitting data:

\begin{enumerate}
\item multinomial maximum posterior;
\item 1-nearest-support distance;
\item rank-32 class-conditional Gaussian;
\item EVM-style Weibull margin approximation;
\item weighted affine support; and
\item sigmoid-calibrated RBF SVM evidence.
\end{enumerate}

A stratified train-only calibration split selects the novelty threshold at
approximately 92\% known coverage. Retention requires at least 90\% known
coverage and 50\% unknown recall on every seed, bounded known-accuracy loss,
competitive review precision, corruption checks, and exact replay.

\section{Persistent discovery}

Rejected examples enter a bounded FIFO buffer of 2,000 records. Five
deterministic windows compare no clustering, streaming centroid
micro-clusters, HDBSCAN, FINCH, and silhouette-selected $k$-means. Groups need
minimum size 10 and at least 70\% purity for scoring. A review ID persists when
consecutive-window member overlap is at least 50\%.

Discovery advancement requires category recovery, review precision, stable
identities, replay, bounded resources, and no semantic publication.

\section{Confirmation and adaptation}

A deterministic oracle simulates delayed human confirmation while preserving
the authority boundary. Outcomes are existing class, new class,
corruption/irrelevant, or unresolved. Operations include no update, native
Gaussian update, rank-16 affine insertion, full class-local refit, and a
full-retraining control.

New-class adaptation uses 50 confirmed supports. Passing requires at least
5 points of target gain, at most 1 point of known regression, at most 2 points
of remaining-unknown loss, stable NLL, graph validity, confirmation lineage,
replay, and rollback.

\section{Routing}

Three independently owned bundles have class orders $\{0,1,2,3\}$,
$\{4,5,6,7\}$, and $\{8\}$. Profiles include centroids and radii, rank-16
Gaussians, compact prototypes, and undercomplete autoencoders. Structural
fingerprints determine compatibility. Empirical profiles shortlist two
bundles.

Authority requires at least 99\% correct-bundle top-1, 99.9\% exhaustive-winner
inclusion and final agreement, 95\% no-confident-route recall, 25\% exact-model
evaluation reduction, lower tail latency, exact update replay, and complete
stale-profile fallback.

\section{Integrated factorial}

Only retained mechanisms enter the integrated study: rank-32 Gaussian
rejection, HDBSCAN or FINCH discovery, delayed confirmation, rank-16 insertion,
and exhaustive routing. A no-grouping review-everything arm is retained as a
burden control. The grouped review budget is 50.

The integrated gate requires at least 75\% of confirmable classes integrated,
25\% review reduction, known and unknown preservation, no unconfirmed
publication, and all transaction contracts.

\chapter{Adaptation-Utility Study}

\section{Episode replay}

The v8 study begins with classes 0--5 and introduces classes 6, 7, and 8 in
order for each of the three seeds, producing nine seed-by-arrival cells. Class
9 remains unknown. Each class is deterministically partitioned into geometry,
anchor, review-candidate, and proxy-validation subsets.

The proxy-validation subset scores a frozen candidate review set but never
becomes adaptation support. Each episode records partition hashes, class order,
selected IDs, confirmation IDs, thresholds, parent and child hashes, utility,
safety metrics, and rollback evidence.

\section{Threshold-transfer experiment}

Adding a class changes score competition. Four rules are compared:

\begin{enumerate}
\item frozen pre-integration threshold;
\item class-count heuristic;
\item global anchor-quantile recalibration; and
\item per-class anchor thresholds under a global fail-closed ceiling.
\end{enumerate}

The experiment consumes no review labels. Stale anchor lineage never lowers
confidence and restores the parent threshold. The selected rule must preserve
known accuracy and remaining-unknown recall over all three seeds.

\section{Sufficient-statistics audit}

The audit measures whether discovery outputs represent the candidate region.
Frozen features include:

\begin{itemize}
\item distance to cluster center;
\item distance to the parent decision boundary;
\item low-rank subspace coverage;
\item covariance trace ratio;
\item omitted-region nearest-neighbor coverage; and
\item class-mode coverage.
\end{itemize}

The audit also ranks measured interface mismatches by failure frequency,
estimated utility, and implementation cost.

\section{Equal-budget review policies}

All eligible arms receive 50 labels:

\begin{description}
\item[Core-selected:] nearest examples to the candidate center.
\item[Boundary-inclusive:] deterministic mixture of core, modes, and
parent-boundary quantiles.
\item[Coverage-selected:] farthest-first $k$-center support.
\item[Utility-selected:] deterministic subsets ranked by a rank-3 proxy adapter
on the disjoint validation partition.
\item[Random-stratified:] seeded sampling over two candidate modes.
\item[Review-everything:] all 300 candidates; diagnostic and ineligible.
\end{description}

The production adapter remains rank 16. All candidates in this controlled
study belong to the arriving class, fixing purity at 100\% and isolating
representativeness.

\section{Statistical analysis}

The primary comparison is paired by seed and arrival class. Mean utility
difference is accompanied by a deterministic 10,000-resample paired percentile
bootstrap interval. Cell signs are reported because a mean can conceal
schedule-specific failures. Advancement requires practical effect, interval,
cell consistency, and safety jointly.

\section{Local residual branch}

Two residual families are preregistered:

\begin{enumerate}
\item local target-class temperature; and
\item bounded class-local affine correction.
\end{enumerate}

The parent freezes the affected region before fitting using selected support,
a bounded nearest-neighbor radius, and component responsibility. Fusion is
gated by this parent-only mask. Passing requires 5 points of utility gain,
99.9\% unaffected-prediction preservation on every seed, known and unknown
safety, local update work, replay, and rollback. At most two attempts are
allowed.

\chapter{Results}

\section{Predictive competitiveness of explicit heads}

The v6.1 weighted rank-32 affine head reaches 91.73\% balanced accuracy,
improving its retained explicit predecessor by 0.97 points. The matched RBF SVM
reaches 96.77\%, leaving a 5.03-point gap and failing the registered 2-point
parity operand. A direct radial family peaks and then deteriorates as component
budget increases, indicating instability rather than monotonic capacity gain.

\begin{table}[ht]
\centering
\caption{Selected predictive results under frozen DINOv2 features.}
\label{tab:predictive}
\begin{tabular}{lrr}
\toprule
Head & Balanced accuracy & Difference from RBF \\
\midrule
Weighted affine rank 32 & 91.73\% & $-5.03$ points\\
Direct radial rank 32 & 90.77\% & $-6.00$ points\\
RBF SVM control & 96.77\% & ---\\
\bottomrule
\end{tabular}
\end{table}

All lifecycle and replay gates pass for the selected explicit head. The result
answers RQ1 negatively for predictive parity but positively for deterministic,
editable artifact behavior. Editability is not evidence of accuracy
leadership.

\section{Acceptance-head results}

The low-rank Gaussian is the only v7 acceptance head to pass every autonomy and
review gate.

\begin{table}[ht]
\centering
\caption{Three-seed mean acceptance-head results.}
\label{tab:acceptance}
\scriptsize
\begin{tabular}{lrrrrc}
\toprule
Head & Coverage & Unknown recall & AUROC & Review precision & Retained\\
\midrule
Max posterior & 91.38\% & 60.17\% & 0.9053 & 69.33\% & No\\
kNN support & 91.83\% & 67.00\% & 0.9045 & 86.67\% & No\\
Low-rank Gaussian & 92.08\% & 86.17\% & 0.9561 & 98.67\% & Yes\\
EVM-style margin & 100.00\% & 0.00\% & 0.5102 & 22.00\% & No\\
Weighted affine & 90.75\% & 53.00\% & 0.8763 & 64.67\% & No\\
RBF evidence & 93.08\% & 12.83\% & 0.4010 & 31.33\% & No\\
\bottomrule
\end{tabular}
\end{table}

This result reinforces the distinction between geometric coverage and
calibrated discrimination. Proper density is retained over raw affine support.

\section{Discovery results}

\begin{table}[ht]
\centering
\caption{Persistent discovery results.}
\label{tab:discovery}
\begin{tabular}{lrrrc}
\toprule
Policy & Group recall & Review precision & Full cells & Pass\\
\midrule
No clustering & 0.00\% & 0.00\% & 0/9 & No\\
Micro-clusters & 0.00\% & 98.67\% & 0/9 & No\\
HDBSCAN & 100.00\% & 100.00\% & 9/9 & Yes\\
FINCH & 83.33\% & 86.39\% & 8/9 & Yes\\
$k$-means & 100.00\% & 98.67\% & 9/9 & Yes\\
\bottomrule
\end{tabular}
\end{table}

HDBSCAN and FINCH produce stable review objects and pass stage gates. The
micro-cluster result shows that sample-level review precision does not imply
separation of distinct unknown classes.

\section{Isolated adaptation and routing}

Confirmed rank-16 insertion raises held-out class-8 success by 42\%, 52\%,
and 36\% across seeds, a 43.33-point mean gain. Known accuracy and
remaining-unknown recall remain within bounds, and all child bundles replay and
roll back exactly.

No routing profile qualifies for authority. Best top-1 is 93.53\%, best winner
inclusion is 99.57\%, and best unknown fallback recall is 71.00\%. Profiles
reduce exact model evaluations by approximately 50\%, but compute savings are
not a safety result. Routing remains shadow-only.

\section{Stage composition failure}

\begin{table}[ht]
\centering
\caption{Integrated v7 delayed-confirmation results.}
\label{tab:integrated}
\begin{tabular}{lrrrr}
\toprule
Discovery & Integrated & Reviews & Review reduction & Unknown recall\\
\midrule
HDBSCAN & 0/3 & 50.00 & 78.78\% & 79.33\%\\
FINCH & 0/3 & 15.33 & 93.49\% & 79.33\%\\
Review everything & 3/3 & 235.67 & 0.00\% & 79.33\%\\
\bottomrule
\end{tabular}
\end{table}

Qualified HDBSCAN and FINCH loops integrate no confirmable classes. Reviewing
all rejections integrates all three, but eliminates the human-budget advantage.
Remaining-unknown recall is 6.83 points below the M39 stage baseline. All
confirmation, graph, replay, rollback, and audit contracts pass. The failure is
therefore one of utility and composition, not transaction correctness.

\section{Threshold transfer}

Frozen, class-count, and global anchor-quantile rules all preserve mean known
accuracy and remaining-unknown recall in the isolated M46 episodes. The global
anchor rule is retained because it directly recalibrates after class-order
change. Per-class thresholds increase unknown rejection but cause 1.875 points
of mean known regression, exceeding the one-point ceiling.

\begin{table}[ht]
\centering
\caption{Threshold transfer over three seeds.}
\label{tab:threshold}
\begin{tabular}{lrrc}
\toprule
Rule & Known drop & Unknown drop & Pass\\
\midrule
Frozen parent & 0.000 & 0.000 & Yes\\
Class-count heuristic & 0.000 & 0.000 & Yes\\
Global anchor quantile & 0.000 & 0.000 & Yes\\
Per-class anchor & 1.875 & $-8.667$ & No\\
\bottomrule
\end{tabular}
\end{table}

\section{Representativeness audit}

\begin{table}[ht]
\centering
\caption{Core and boundary-inclusive support diagnostics.}
\label{tab:coverage}
\small
\begin{tabular}{crrrr}
\toprule
Seed & Core NN cov. & Boundary NN cov. & Core trace ratio & Boundary trace ratio\\
\midrule
11 & 34.04\% & 99.68\% & 0.663 & 1.078\\
23 & 35.47\% & 99.78\% & 0.658 & 1.096\\
37 & 39.14\% & 99.37\% & 0.679 & 1.068\\
\bottomrule
\end{tabular}
\end{table}

Core and boundary sets both cover two coarse modes, so mode count alone misses
the defect. Nearest-neighbor support and covariance reveal that core sets omit
most of the rejected region. Boundary-inclusive sets nearly recover full
support, although the cheap M46 proxy utility improves by only 0.259 points.

\section{Equal-budget utility}

\begin{table}[ht]
\centering
\caption{Mean v8 adaptation utility over nine cells.}
\label{tab:utility}
\begin{tabular}{lrrr}
\toprule
Review policy & Mean $\utility_e$ & Labels & Remaining-unknown recall\\
\midrule
Core-selected & 6.36 & 50 & 59.72\%\\
Boundary-inclusive & 9.71 & 50 & 75.98\%\\
Coverage-selected & \textbf{10.62} & 50 & \textbf{82.50\%}\\
Utility-selected & 9.96 & 50 & 78.15\%\\
Random-stratified & 9.70 & 50 & 75.41\%\\
Review-everything & 11.26 & 300 & 86.13\%\\
\bottomrule
\end{tabular}
\end{table}

Utility selection beats core selection in all nine paired cells. The mean
difference is 3.593 points with paired 95\% bootstrap interval
[2.794, 4.382]. The sign and interval operands pass, but the 5-point practical
operand fails. Six cells also exceed the 2-point remaining-unknown-recall band.

\begin{table}[ht]
\centering
\caption{Pivotal preregistered gate.}
\label{tab:pivotal}
\begin{tabular}{lrc}
\toprule
Operand & Observation & Pass\\
\midrule
Mean gain over core $\geq 5.0$ & 3.593 & No\\
Paired 95\% lower bound $>0$ & 2.794 & Yes\\
Positive cells $\geq 7/9$ & 9/9 & Yes\\
Budget, safety, transaction conjunction & 6 recall violations & No\\
\bottomrule
\end{tabular}
\end{table}

Coverage selection is the strongest eligible deterministic arm, slightly
outperforming utility selection. The rank-3 proxy therefore does not perfectly
rank rank-16 production outcomes. Review-everything achieves the highest raw
utility but is ineligible at six times the budget.

\section{Local residuals}

Both residual attempts preserve 100\% of unaffected predictions, improving on
the earlier A3 minimum of 84.2\%. The local temperature adds 0.000 points of
mean utility. The bounded affine correction adds 0.204 points. Maximum known
regression is 0.125 points, but maximum remaining-unknown-recall drop is 4.333
points. Neither residual is retained.

This result separates locality from utility. Explicit parent-scoped fusion
repairs the scope bug, but a correctly localized update is not necessarily a
useful or safe update.

\section{Final disposition}

M47 produces registered Outcome D. Learned selection, same-harness head
continuation, and untouched M50 lifecycle qualification are blocked. The
positive bootstrap interval is reported as a diagnostic finding, not promoted
to primary success.

\chapter{Discussion}

\section{Answer to RQ1}

Compact explicit geometry did not reach predictive parity with the matched RBF
control. The weighted affine head remained 5.03 points behind. Proper
class-conditional Gaussian density also outperformed weighted affine support
for rejection. These outcomes do not remove the governance value of explicit
artifacts, but they reject the premise that editability itself implies
competitive discrimination.

\section{Answer to RQ2}

Independently qualified stages did not compose. HDBSCAN achieved perfect
stage-wise group metrics, and confirmed insertion delivered a large isolated
gain, yet their grouped loop integrated no class. The review-everything control
showed that the adapter could succeed when given broader support. The failure
therefore lies at the clusterer-to-review-to-adapter interface.

\section{Answer to RQ3}

Selecting for downstream utility improves over core selection consistently:
9/9 cells and a positive bootstrap interval. The improvement is nevertheless
only 3.593 points and violates remaining-unknown safety in six cells. The
answer is therefore qualified: adaptation utility is a more discriminating
endpoint and reveals a real representativeness signal, but the implemented
utility selector does not qualify the lifecycle.

\section{Answer to RQ4}

Transactional controls successfully distinguish predictive failure from
systems failure. No unconfirmed class is published, every rollback restores
the exact parent, stale lineages fail closed, and residual scope can be enforced
to preserve all unaffected predictions. These controls do not create accuracy,
but they make negative results attributable and reproducible.

\section{Pure cores as insufficient statistics}

Density-based discovery naturally prefers high-cohesion regions. Such cores
are attractive review objects because they reduce ambiguity and label noise.
The adapter, however, estimates a class region. It needs variance directions,
boundary examples, and low-density modes. The measured 34--39\% support
coverage provides a direct explanation for the v7 failure.

This distinction generalizes beyond the specific adapter. A prototype learner,
density estimator, margin model, or discriminative head can require different
support statistics. Review selection should therefore be conditioned on the
downstream update mechanism.

\section{Adaptation and unknown preservation conflict}

Broader support makes the new class easier to recognize, but it also expands
acceptance into regions occupied by future unknowns. This creates a structural
tradeoff. Utility on known plus discovered classes and recall on remaining
unknowns cannot be optimized independently.

A future objective should be explicitly constrained:

\begin{equation}
\max_{\pi}\ \mathbb{E}[\utility_e(\pi)]
\quad\text{subject to}\quad
\Delta\ba_{\text{known}}\geq-\epsilon,\quad
\Delta R_{\text{unknown}}\geq-\delta,\quad
\operatorname{cost}(\pi)\leq B.
\end{equation}

Here $\pi$ is the review policy. Unknown preservation is not a post-hoc
diagnostic; it is part of the optimization problem.

\section{Why coverage beats the utility proxy}

The utility selector evaluates deterministic subsets with a rank-3 proxy, then
deploys a rank-16 production adapter. This mismatch introduces proxy regret.
Farthest-first coverage has no such model mismatch and reaches 10.62 points,
compared with 9.96 for utility selection.

Possible improvements include a better-fidelity cheap adapter, uncertainty in
the subset ranking, direct constrained selection, or a learned policy trained
across episodes. The registered kill switch correctly blocks these
continuations after the pivotal endpoint fails; they remain future hypotheses.

\section{Methodological implications}

The results support four methodological recommendations:

\begin{enumerate}
\item Register a system-level endpoint before optimizing stages.
\item Audit producer outputs against consumer sufficient statistics.
\item Treat unknown recall as a lifecycle state variable after every update.
\item Separate transaction correctness, locality, and predictive utility.
\end{enumerate}

\chapter{Validity, Limitations, and Ethics}

\section{Internal validity}

The use of frozen partitions, seeds, parent hashes, and deterministic replay
reduces accidental variation. Paired cells isolate policy differences. The
validation partition is disjoint from adaptation support. Fail-closed tests
guard common lineage and leakage errors.

Internal validity remains limited by the deterministic oracle and controlled
candidate purity. Real review candidates may mix classes, and real reviewers
may disagree. The utility proxy and production adapter differ in rank.

\section{Construct validity}

Balanced accuracy gives equal weight to classes and is appropriate for the
episode inventory, but it does not capture calibration, severity-weighted
errors, or real deployment value. The fixed 50-label budget approximates human
cost but does not model variable review difficulty.

Nearest-neighbor support coverage and covariance ratios diagnose
representativeness but are not complete sufficient statistics. Their role is
explanatory rather than universal.

\section{External validity}

The realized study uses CIFAR-10 and frozen DINOv2-small features. Results may
change with more classes, long-tailed data, domain shift, natural streams,
larger encoders, or representation adaptation. The held-out class order is
specific. The study does not establish performance for medical, legal,
security, or other high-stakes applications.

\section{Statistical conclusion validity}

There are nine paired development cells. The bootstrap interval describes this
bounded schedule, not a broad population of domains. The untouched M50 schedule
remains sealed, so there is no independent final confirmation. The
preregistered practical threshold prevents a positive interval from being
overinterpreted.

\section{Ethics and governance}

Open-world systems can assign meaning to coherent but irrelevant, private, or
harmful groups. Automated class creation may amplify surveillance,
discrimination, or feedback loops. This work prohibits autonomous semantic
publication and records review authority, provenance, and rollback.

Deployment would additionally require privacy-aware retention, deletion and
appeal procedures, reviewer disagreement models, subgroup audits, access
control, security review, and monitoring for distributional feedback. The
experiments use public benchmark representations and simulated review; they do
not involve human subjects or sensitive personal data.

\section{Claim boundary}

The evidence supports the claims that:

\begin{itemize}
\item adaptation utility is more discriminating than stage purity;
\item boundary- and coverage-aware review improves utility over density cores
in the frozen harness;
\item global anchor recalibration transfers safely in the isolated M46
episodes; and
\item parent-only fusion repairs the observed locality leakage.
\end{itemize}

It does not support claims of closed-set parity, SDF necessity, autonomous
class creation, learned-selector superiority, authoritative routing,
head-substitution superiority, untouched confirmation, or qualified
end-to-end operation.

\chapter{Post-Lifecycle Geometric Support Studies}

\section{Motivation and methodological evolution}

The negative lifecycle result raised a narrower representation question:
whether frozen class features occupy volumes, lie near class boundaries, or
concentrate around higher-codimension supports. V9 first tested shell occupancy
around the frozen A2 zero-level sets. None of 48 class-by-seed-by-score
diagnostics contained enough negative own-class interior to form the registered
equal-mass comparison, closing only that specific shell hypothesis.

The distinct v9 affine-tube branch produced predictive and unknown-recall
signal at ranks 16 and 32, but accepted 100\% of 8x tangent probes. This failure
identified a score-unit defect: orthogonal residuals and tangent overshoot were
not calibrated on comparable scales. V10 therefore registered dimensionless
residuals, rank-normalized tangent penalties, calibration-only penalty
selection, source-versus-system safety evidence, and controlled identifiability
before any real-feature screen.

\section{Controlled and real-feature results}

In ambient dimension 64, the implementation recovered true ranks 8, 16, and 32
in all nine seed cells, preserved 91.35\% pooled in-support coverage, and
rejected every 8x tangent probe. Two-patch atlases improved controlled curved
and two-mode supports, while full-dimensional volume, spherical-shell,
random-orientation, and random-label controls behaved as registered.

The frozen seed-11 screen then evaluated 18 rank-by-extent-by-scale cells.
Sixteen were calibration-infeasible because minimum-system acceptance remained
too high under tangent extrapolation. Only rank 8 with 0.99 extents was
feasible. One scale policy regressed known balanced accuracy by 5.75 points;
the other gained 0.875 points, missing the registered 1-point threshold. No
global tube was retained, the real-data condition for a local atlas did not
occur, and lifecycle testing remained blocked.

\section{Resulting claim boundary}

The combined v9--v10 record does not prove that frozen features fill volumes,
nor does it universally disprove manifold structure. It shows that the tested
zero-level shells were not calibrated supports and that the registered bounded
affine family, despite controlled identifiability, did not satisfy joint
predictive and open-space gates on the frozen features. Nonlinear manifolds or
different learned representations remain untested rather than supported.

\chapter{Reproducibility}

\section{Artifact structure}

The repository stores implementation code, tests, plans, claim ledgers,
machine-readable evidence, and immutable artifact indexes. Key locations are:

\begin{itemize}
\item \texttt{analysis/V8\_FINAL\_CLAIM\_LEDGER.md};
\item \texttt{analysis/RESEARCH\_IMPLEMENTATION\_PLAN\_v8.md};
\item \texttt{logs/results/v8/m46\_diagnostics/};
\item \texttt{logs/results/v8/m47\_review\_utility/};
\item \texttt{logs/results/v8/m49\_local\_residuals/}; and
\item \texttt{logs/results/v8/final\_outcome\_replay/};
\item \texttt{analysis/V10\_FINAL\_CLAIM\_LEDGER.md}; and
\item \texttt{logs/results/v10/m62\_final\_replay/}.
\end{itemize}

\section{Final replay}

From the repository root on Windows:

\begin{verbatim}
.\.venv\Scripts\python.exe -m experiments.tier1.eval_v8_final_replay
\end{verbatim}

The command verifies four immutable milestone indexes, 11 indexed artifacts,
and six conclusion operands. It runs twice and requires byte-identical output,
loads no training data, and opens no final labels. The complete repository gate
contains 374 tests at finalization.

\section{Building this report}

From the \texttt{analysis} directory:

\begin{verbatim}
tectonic MS_THESIS_REPORT.tex
\end{verbatim}

Tectonic resolves standard LaTeX packages automatically. The source uses an
inline bibliography and no external figures, reducing archival dependencies.
Institution-specific title, approval, declaration, margins, and binding rules
must be adapted before university submission.

\chapter{Conclusions and Future Work}

\section{Conclusions}

This thesis investigated open-world recognition as a composed, transactional
lifecycle. The program first tested whether explicit geometric artifacts could
serve as competitive class heads. They replayed and rolled back reliably but
remained behind discriminative controls. It then qualified rejection,
discovery, and confirmed adaptation stages independently. Those stages failed
when composed.

The successor study replaced stage metrics with adaptation utility at a fixed
review budget. It found that density cores omit most of the candidate region
and that representative selection improves downstream utility. The positive
effect is consistent but insufficient: it misses the registered practical
threshold and damages remaining-unknown recall.

The principal conclusion is methodological. Open-world learning is not a
pipeline whose stages can be optimized independently. Every update changes the
known inventory, acceptance threshold, future unknown space, and routing
state. Qualification must therefore occur at the lifecycle level under joint
utility, safety, cost, lineage, locality, and rollback constraints.

The post-lifecycle geometric studies reached the same methodological boundary
from another direction. Synthetic identifiability and correctly scaled safety
scores did not guarantee a useful real-feature support model. Most affine cells
failed calibration safety, and the best safe cell missed its predictive gate.

\section{Future work}

The next research program should begin with new registration rather than reopen
Outcome D. Promising directions are:

\begin{enumerate}
\item constrained review selection that predicts both adaptation gain and
remaining-unknown loss;
\item higher-fidelity but bounded proxy adapters;
\item leave-one-episode-out selection learners with explicit uncertainty;
\item evaluation over larger category inventories and natural class-arrival
streams;
\item real human-review studies with disagreement and variable cost;
\item representation adaptation under strict parent and rollback contracts;
\item mechanism-matched comparisons among Gaussian, EVM, kNN, and explicit
geometry under the same lifecycle endpoint; and
\item router training that treats unknown fallback as a primary safety target.
\end{enumerate}

Any continuation should retain the central discipline of this work: practical
effect thresholds and safety constraints must be registered before tuning, and
failed branches must remain visible.

\begin{thebibliography}{32}

\bibitem{bendale2015}
A. Bendale and T. Boult.
\newblock Towards open world recognition.
\newblock In \emph{CVPR}, 2015.

\bibitem{bendale2016}
A. Bendale and T. Boult.
\newblock Towards open set deep networks.
\newblock In \emph{CVPR}, 2016.

\bibitem{rudd2018}
E. Rudd, L. Jain, W. Scheirer, and T. Boult.
\newblock The extreme value machine.
\newblock \emph{IEEE TPAMI}, 2018.

\bibitem{masud2011}
M. Masud et al.
\newblock Classification and novel class detection in concept-drifting data
streams under time constraints.
\newblock \emph{IEEE TKDE}, 2011.

\bibitem{haque2016sand}
A. Haque et al.
\newblock SAND: Semi-supervised adaptive novel class detection and
classification over data stream.
\newblock In \emph{AAAI}, 2016.

\bibitem{haque2016echo}
A. Haque et al.
\newblock ECHO: A data stream classification method for detecting novel class
drift.
\newblock In \emph{ICDE}, 2016.

\bibitem{faria2016}
E. de Faria, A. Gama, and A. Carvalho.
\newblock Novelty detection algorithm for data streams multi-class problems.
\newblock \emph{Data Mining and Knowledge Discovery}, 2016.

\bibitem{mu2017}
X. Mu, F. Zhu, J. Du, E.-P. Lim, and Z.-H. Zhou.
\newblock Classification under streaming emerging new classes.
\newblock \emph{IEEE TKDE}, 2017.

\bibitem{vaze2022}
S. Vaze et al.
\newblock Generalized category discovery.
\newblock In \emph{CVPR}, 2022.

\bibitem{cao2022}
K. Cao, M. Brbi\'c, and J. Leskovec.
\newblock Open-world semi-supervised learning.
\newblock In \emph{ICLR}, 2022.

\bibitem{zhang2022}
Z. Zhang et al.
\newblock Grow and Merge: A unified framework for continuous categories
discovery.
\newblock In \emph{NeurIPS}, 2022.

\bibitem{zhao2023}
B. Zhao and O. Mac Aodha.
\newblock Incremental generalized category discovery.
\newblock In \emph{ICCV}, 2023.

\bibitem{aljundi2017}
R. Aljundi, P. Chakravarty, and T. Tuytelaars.
\newblock Expert Gate: Lifelong learning with a network of experts.
\newblock In \emph{CVPR}, 2017.

\bibitem{gururangan2021}
S. Gururangan et al.
\newblock DEMix layers: Disentangling domains for modular language modeling.
\newblock arXiv:2108.05036, 2021.

\bibitem{li2022}
M. Li et al.
\newblock Branch-Train-Merge: Embarrassingly parallel training of expert
language models.
\newblock arXiv:2208.03306, 2022.

\bibitem{hendrycks2017}
D. Hendrycks and K. Gimpel.
\newblock A baseline for detecting misclassified and out-of-distribution
examples in neural networks.
\newblock In \emph{ICLR}, 2017.

\bibitem{lee2018}
K. Lee et al.
\newblock A simple unified framework for detecting out-of-distribution samples
and adversarial attacks.
\newblock In \emph{NeurIPS}, 2018.

\bibitem{mukhoti2023}
J. Mukhoti et al.
\newblock Deep deterministic uncertainty: A new simple baseline.
\newblock \emph{TMLR}, 2023.

\bibitem{sun2022}
Y. Sun et al.
\newblock Out-of-distribution detection with deep nearest neighbors.
\newblock In \emph{ICML}, 2022.

\bibitem{vanamersfoort2020}
J. van Amersfoort et al.
\newblock Uncertainty estimation using a single deep deterministic neural
network.
\newblock In \emph{ICML}, 2020.

\bibitem{liu2020}
J. Liu et al.
\newblock Simple and principled uncertainty estimation with deterministic deep
learning via distance awareness.
\newblock In \emph{NeurIPS}, 2020.

\bibitem{ester1996}
M. Ester, H.-P. Kriegel, J. Sander, and X. Xu.
\newblock A density-based algorithm for discovering clusters in large spatial
databases with noise.
\newblock In \emph{KDD}, 1996.

\bibitem{campello2013}
R. Campello, D. Moulavi, and J. Sander.
\newblock Density-based clustering based on hierarchical density estimates.
\newblock In \emph{PAKDD}, 2013.

\bibitem{sarfraz2019}
S. Sarfraz, V. Sharma, and R. Stiefelhagen.
\newblock Efficient parameter-free clustering using first neighbor relations.
\newblock In \emph{CVPR}, 2019.

\bibitem{requicha1980}
A. Requicha.
\newblock Representations for rigid solids: Theory, methods, and systems.
\newblock \emph{ACM Computing Surveys}, 1980.

\bibitem{park2019}
J. Park et al.
\newblock DeepSDF: Learning continuous signed distance functions for shape
representation.
\newblock In \emph{CVPR}, 2019.

\bibitem{tipping1999}
M. Tipping and C. Bishop.
\newblock Probabilistic principal component analysis.
\newblock \emph{JRSS B}, 1999.

\bibitem{ghahramani1996}
Z. Ghahramani and G. Hinton.
\newblock The EM algorithm for mixtures of factor analyzers.
\newblock Technical Report CRG-TR-96-1, University of Toronto, 1996.

\bibitem{oquab2024}
M. Oquab et al.
\newblock DINOv2: Learning robust visual features without supervision.
\newblock \emph{TMLR}, 2024.

\end{thebibliography}

\appendix

\chapter{Registered Gates}

\begin{longtable}{
>{\raggedright\arraybackslash}p{1.1in}
>{\raggedright\arraybackslash}p{4.4in}}
\toprule
\textbf{Stage} & \textbf{Principal advancement conditions}\\
\midrule
\endhead
Predictive head & Within 2 points of RBF control, lifecycle checks, exact replay
and rollback.\\
Rejection & Every seed: at least 90\% known coverage, 50\% unknown recall,
bounded known loss, competitive review precision, exact replay.\\
Discovery & Category recovery, at least 70\% precision, stable review IDs,
resource limits, no semantic publication.\\
Adaptation & At least 5 points target gain, at most 1 point known drop, at most
2 points unknown-recall drop, confirmation, replay, rollback.\\
Routing & At least 99\% top-1, 99.9\% winner inclusion and agreement, 95\%
unknown fallback, and 25\% compute reduction.\\
Integrated v7 & At least 75\% classes integrated, 25\% review reduction,
known/unknown preservation, all safety contracts.\\
Utility v8 & At least 5 points over core, positive paired interval, at least
7/9 positive cells, budget and safety conjunction.\\
Residual & At least 5 points utility, 99.9\% unaffected preservation per seed,
known/unknown safety, local work, replay, rollback.\\
\bottomrule
\end{longtable}

\chapter{Final Branch Disposition}

\begin{longtable}{
>{\raggedright\arraybackslash}p{1.1in}
>{\raggedright\arraybackslash}p{0.9in}
>{\raggedright\arraybackslash}p{3.4in}}
\toprule
\textbf{Branch} & \textbf{Status} & \textbf{Reason}\\
\midrule
\endhead
Weighted affine parity & Failed & 91.73\%, 5.03 points behind RBF.\\
Gaussian rejection & Passed & Only head satisfying all v7 acceptance gates.\\
HDBSCAN discovery & Passed stage & Perfect stage-wise group recall and
precision.\\
FINCH discovery & Passed stage & Mechanistically distinct retained control.\\
New-class insertion & Passed stage & 43.33-point isolated target gain with
rollback.\\
Authoritative routing & Closed & Accuracy and unknown-fallback gates missed.\\
Integrated grouped loop & Failed & 0/3 classes integrated.\\
M46 threshold transfer & Passed & Global anchor quantile retained.\\
M47 utility selection & Outcome D & 3.593 points below 5-point gate; six safety
violations.\\
M48 learned selector & Blocked & Binding M47 kill switch.\\
M49 residuals & Closed & Locality passed; utility and unknown safety failed.\\
M50/E12 confirmation & Blocked & Untouched schedule remains sealed.\\
\bottomrule
\end{longtable}

\chapter{Artifact-Only Verification}

The final v8 verifier reports:

\begin{itemize}
\item four verified milestone indexes;
\item 11 verified indexed artifacts;
\item six reproduced conclusion operands;
\item byte-identical output across two runs;
\item no training-data loading; and
\item no final-label access.
\end{itemize}

The final claim ledger is
\texttt{analysis/V8\_FINAL\_CLAIM\_LEDGER.md}. The paper-length presentation is
\texttt{analysis/FINAL\_RESEARCH\_PAPER.tex}. This thesis report provides the
expanded academic narrative; it does not introduce new experimental claims.

\end{document}
